\begin{center}\begin{tikzpicture}[font=\small,>=stealth,line join=round]
\draw[gray!40,rounded corners=3pt] (0,0) rectangle (14.2,8.3);
\node[text width=12.6cm,align=center] at (7.1,7.65) {一般水平轴 $y=3$：横坐标相同，高度之和为 $6$};
\begin{scope}[shift={(1.66,2.69)},x=1.3199999999999998cm,y=0.47000000000000003cm]
\draw[->,gray] (-0.5,0)--(4.5,0) node[right]{$x$};\draw[->,gray] (0,-2)--(0,8) node[above]{$y$};
\draw[violet,dashed] (-.3,3)--(4.2,3) node[right]{$y=3$};\draw[gray,dashed] (2,1)--(2,5);
\begin{scope}\clip (-0.5,-2) rectangle (4.5,8);
\draw[blue,thick,domain=-0.5:4.5,samples=180] plot(\x,{(\x-2)^2+1});
\draw[orange!85!black,thick,domain=-0.5:4.5,samples=180] plot(\x,{5-(\x-2)^2});
\end{scope}
\fill[blue] (2,1) circle(2.2pt);\node[below right=5pt,blue] at (2,1) {$P$};
\fill[orange!85!black] (2,5) circle(2.2pt);\node[above right=5pt,orange!85!black] at (2,5) {$Q$};
\end{scope}
\node[anchor=west,align=left,text width=4.8cm] at (8.65,4.5) {蓝：$f(x)=(x-2)^2+1$\\[8pt]橙：$g(x)=5-(x-2)^2$\\[8pt]$P(2,1)\to Q(2,5)$\\[8pt]$g(x)=6-f(x)$};
\node[text width=12.6cm,align=center] at (7.1,.65) {一般形式：关于 $y=b$ 相互对称的另一个函数为 $g(x)=2b-f(x)$。};
\end{tikzpicture}\end{center}
